\documentclass[10pt,a4paper]{article}
\usepackage{amsmath,amssymb,amsthm,graphicx}
\title{Minimal vanishing spaces of random sets and high-rate Reed-Muller resilience}
\author{Siddharth Bhandari, Prahladh Harsha, Ramprasad Saptharishi, Srikanth Srinivasan}
\date{doi:10.4230/LIPIcs.CCC.2022.31; CC BY4.0}
\begin{document}
\maketitle
\noindent\textit{Editorial note.} Original authors: Siddharth Bhandari, Prahladh Harsha, Ramprasad Saptharishi, Srikanth Srinivasan. Official CCC2022 article31, DOI10.4230/LIPIcs.CCC.2022.31.\par\medskip
\noindent\textit{Editorial note.} License: https://creativecommons.org/licenses/by/4.0/\par\medskip
\noindent\textit{Editorial note.} Faithful editable TeX transcription of the complete formal author Theorem1.1 proof, Lemma3.4 and Claims3.5/3.6 continuations, required definitions and cited weight-distribution prerequisites, and direct Reed-Muller consequences with all author Section4 derivations and complete bibliography. Official CC BY4.0 version only; the arXiv BY-NC-ND4.0 version is identity/license evidence and its content is not used. Original published expressions and conditions retained without mathematical correction or additional proofs.\par\medskip
\noindent\textit{Allowed standard prerequisites (editorial).}\begin{itemize}
\item Sberlo--Shpilka weight distribution bounds. Exact author Theorems2.1/2.2[13] and parameter footnote retained.
\item Reed-Muller erasure-to-error decoding reduction. Exact cited author Theorem4.6[1,12] retained.
\end{itemize}\medskip
\section*{Author notation and preliminaries}
We use $[m]$ to represent $\{1,2,3,\ldots,m\}$ and $\binom m{\le r}$ to denote $\binom m0+\cdots+\binom mr$. We denote by $\binom{[m]}r$ and $\binom{[m]}{\le r}$ the sets $\{S\subseteq[m]:|S|=r\}$ and $\{S\subseteq[m]:|S|\le r\}$ respectively.
We shall abuse notation and use $\mathrm{RM}(m,r)$ to denote the vector space of all $m$-variate multivariate polynomials in $\mathbb F_2[x_1,\ldots,x_m]$ of degree at most $r$. Throughout the paper, the parameter $m$ will be unchanged and we will avoid mentioning it for brevity. Let $n=2^m$.
For parameters $m,r$, define the matrix $E(m,r)$ as the $\binom m{\le r}\times2^m$-matrix whose columns are indexed by elements in $\mathbb F_2^m$ and rows by multlinear $m$-variate monomials of degree at most $r$ and whose $z$-th column is the vector consisting of the evaluation of all multilinear $m$-variate degree $\le r$ monomials on $z$.
For a polynomial $P\in\mathrm{RM}(m,r)$, we use $\mathrm{wt}(P)$ to denote the fractional Hamming weight:
\[\mathrm{wt}(P):=\frac{|\{z\in\mathbb F_2^m:P(z)\ne0\}|}{2^m}.\]
We shall say that a polynomial $P\in\mathrm{RM}(m,r)$ is $\eta$-biased if $|\mathrm{wt}(P)-1/2|\ge\eta$. We shall also use $\mathrm{RM}_\eta(m,r)$ to denote the set of all $\eta$-biased polynomials in $\mathrm{RM}(m,r)$. For a real number $\alpha\in(0,1)$, we denote by $\mathrm{WtDist}_{m,r}(\alpha)$ the number of polynomials of (fractional) weight at most $\alpha$ in $\mathrm{RM}(m,r)$, i.e., $\mathrm{WtDist}_{m,r}(\alpha):=|\{P\in\mathrm{RM}(m,r):\mathrm{wt}(p)\le\alpha\}|$. All complexity notations used in the paper are with respect to $m$ as the growing parameter.
\paragraph{Theorem 2.1 (Sberlo and Shpilka [13]: bounds for low-weight codewords).} For any $\ell\ge1$, we have
\[\mathrm{WtDist}_{m,r}(2^{-\ell})\le\operatorname{exp}_2\left(O(m^4)+17(c_\gamma\ell+d_\gamma)\gamma^{\ell-1}\binom m{\le r}\right)\]
where $c_\gamma=1/(1-\gamma)$ and $d_\gamma=(2-\gamma)/(1-\gamma)^2$. In particular, if $\gamma\le1/2$, we have
\[\mathrm{WtDist}_{m,r}(2^{-\ell})\le\operatorname{exp}_2\left(O(m^4)+O(\ell\gamma^{\ell-1})\binom m{\le r}\right).\tag{1}\]
\paragraph{Theorem 2.2 (Sberlo and Shpilka [13]: bounds for medium-weight codewords).} Assume that $\gamma\in(0,(1/2)-\Omega(1))$. For a positive integer $\ell$ such that $\ell/m$ upper bounded by a small enough constant,\footnote{[13, Theorem 1.3] only guarantees this for $\ell=o(m)$ but [13, Remark 1.1] after the theorem statement says that in this setting it holds for $\ell=\Omega(m)$.} we have
\[\mathrm{WtDist}_{m,r}\left(\frac12-2^{-\ell}\right)\le\operatorname{exp}_2\left(O(m^4)+(1-2^{-c(\gamma,\ell)})\binom m{\le r}\right),\tag{2}\]
where $c(\gamma,\ell)=O(\max\{\gamma^2\ell,\gamma\})$.
\section*{Degree-$r$ vanishing spaces and closures}
We first define the notion of a vanishing space. This is similar to the notion of vanishing ideals in basic algebraic geometry but we refer to them as vanishing space instead to stress that we are studying them as a vector space and not an ideal.
\paragraph{Definition 3.1 (degree-$r$ vanishing spaces).} For a set $Z\subseteq\mathbb F_2^m$, we use $I_r(Z)$ to denote the degree-$r$ vanishing space defined as
\[I_r(Z):=\{P\in\mathrm{RM}(m,r):P(z)=0\text{ for all }z\in Z\}.\]
Related to the vanishing ideals is also the notion of the degree-$r$ closure (similar to the Zariski closure in standard algebraic geometry but restricted to the setting of $\mathbb F_2^m$), which is the set of points on which every polynomial in the degree-$r$ vanishing space vanishes.
\paragraph{Definition 3.2 (degree-$r$ closure).} For a set $Z\subseteq\mathbb F_2^m$, we use $\mathrm{Closure}_r(Z)$ to denote the degree-$r$ closure of $Z$ defined as
\[\mathrm{Closure}_r(Z):=\{u\in\mathbb F_2^m:P(u)=0\text{ for all }P\in I_r(Z)\}.\]
The above notion can be equivalently defined as the set of all $u\in\mathbb F_2^m$ such that column of $E(m,r)$ indexed by $u$ is in the span of columns of $E(m,r)$ indexed by $z\in Z$. That is,
\[\mathrm{Closure}_r(Z)=\{u\in\mathbb F_2^m:u^{(r)}\in\operatorname{span}\{z^{(r)}:z\in Z\}\},\]
where $z^{(r)}$ denotes the $\binom m{\le r}$-dimension vector of evaluations of all $m$-variate degree at most $r$ monomials at the point $z$. For any set $Z\subseteq\mathbb F_2^m$ of size $k$, note that $I_r(Z)$ is a vector space of dimension at least $\binom m{\le r}-k$, as each constraint $P(z)=0$ adds one homogeneous linear constraint on the ambient space $\mathrm{RM}(m,r)$. In fact, it is easy to see that $I_r(Z)$ has rank $\binom m{\le r}-k$ if and only if each point of $Z$ is not in the closure of the previous points.
\paragraph{Observation 3.3 (vanishing ideals of minimal rank).} Let $Z=\{z_1,\ldots,z_k\}\subseteq\mathbb F_2^m$. Then,
\[\dim I_r(Z)=\binom m{\le r}-k\iff(\forall i=1,\ldots,k-1:z_i\notin\mathrm{Closure}_r(\{z_1,\ldots,z_{i-1}\})).\]

\section*{Complete author main result and direct coding consequences}
\paragraph{Theorem 1.1 (dimension of degree-$r$ vanishing space of random set).} There exists a constant $\gamma_0>0$ such that for all $\epsilon>0$ the following is true. Let $K=(1-\epsilon)\binom m{\le r}$ and $r<\gamma_0m$. Then,
\[\Pr_{Z\subseteq\mathbb F_2^m}\left[\dim(I_r(Z))=\binom m{\le r}-K\right]=1-o(1)\]
where $Z=\{z_1,\ldots,z_K\}$ is a uniformly random set of $K$ distinct points in $\mathbb F_2^m$. In other words, the probability that the set of columns $\{z_1^{(r)},z_2^{(r)},\ldots,z_K^{(r)}\}$ is not linearly independent is $o(1)$, where $z^{(r)}$ denotes the vector consisting of evaluations of all multilinear $m$-variate degree $\le r$ monomials at the point $z$.
\paragraph{Theorem 1.2 (high-degree RM codes under erasures).} There exists a constant $\gamma_0>0$ such that for all $\epsilon>0$ the following is true. Let $m,d$ be growing parameters with $d>m(1-\gamma_0)$. Then, the code $\mathrm{RM}(m,d)$ can correct $K=(1-\epsilon)\binom m{\le m-d-1}$ random errors with probability $1-o(1)$.
\paragraph{Corollary 1.3 (high-degree RM codes under errors).} There exists a constant $\gamma_0>0$ such that for all $\epsilon>0$ the following is true. Let $m,r$ be growing parameters with $r<\gamma_0m$. Then, the code $\mathrm{RM}(m,m-2r-2)$ can be efficiently decode from $K=(1-\epsilon)\binom m{\le r}$ random errors.

\section*{Author Section 3.1: dimension of vanishing spaces}
In this section, we will be interested in studying the vanishing spaces and closures of a random set $Z$ obtained by picking $K=(1-\epsilon)\binom m{\le r}$ points from $\mathbb F_2^m$ (where $\epsilon>0$ is some constant). We will use the following lemma to prove the above theorem.
\paragraph{Lemma 3.4 (size of degree $r$-vanishing space of random set).} There exists a constant $\gamma_0>0$ such that for all $\epsilon>0$ the following is true. Let $K=(1-\epsilon)\binom m{\le r}$ and $r<\gamma_0m$. Then,
\[\mathbb E\left[\frac{|I_r(Z)|}{\operatorname{exp}_2\left(\binom m{\le r}\right)}\right]=2^{-K}(1+o(1))\]
where $Z=\{z_1,\ldots,z_K\}$ is a uniformly random set of $K$ distinct points in $\mathbb F_2^m$. We first use the above lemma to prove Theorem 1.1. Lemma 3.4 is proved in Subsection 3.2.
\paragraph{Proof of Theorem 1.1.} Let $p$ be the probability that $\dim(I_r(Z))=\binom m{\le r}-K$. Note that $\dim(I_r(Z))$ is always at least $\binom m{\le r}-K$, as the space $I_r(Z)$ is a subspace of $\mathrm{RM}(m,r)$ defined by $K$ linear equations. Thus,
\begin{align*}\mathbb E\left[\frac{|I_r(Z)|}{\operatorname{exp}_2\left(\binom m{\le r}\right)}\right]&=\mathbb E\left[\operatorname{exp}_2\left(\dim(I_r(Z))-\binom m{\le r}\right)\right]\\&\ge p\operatorname{exp}_2(-K)+(1-p)2\operatorname{exp}_2(-K).\end{align*}
Combining this with Lemma 3.4 we get that
\[p\operatorname{exp}_2(-K)+(1-p)2\operatorname{exp}_2(-K)\le2^{-K}(1+o(1))\]
which gives $1-p\le o(1)$. Hence, we get that $p=1-o(1)$.\qed
\section*{Author Section 3.2: proof of Lemma 3.4}
We start with reformulating the problem of bounding size of $I_r(Z)$ to computing a certain weighted sum of the polynomials in $\mathrm{RM}(m,r)$. Let $Q$ be a uniformly random polynomial chosen from $\mathrm{RM}(m,r)$. Then,
\begin{align*}
\mathbb E\left[\frac{|I_r(Z)|}{\operatorname{exp}_2\left(\binom m{\le r}\right)}\right]&=\mathbb E_{Z,Q}[\mathbf1[Q\in I_r(Z)]]\\&=\mathbb E_Q[\mathbb E_Z[\mathbf1[Q\in I_r(Z)]]]\\
&=\sum_{P\in\mathrm{RM}(m,r)}\operatorname{exp}_2\left(-\binom m{\le r}\right)\frac{\binom{(1-\mathrm{wt}(P))2^m}K}{\binom{2^m}K}\\
&=\operatorname{exp}_2(-K)\left[\sum_{P\in\mathrm{RM}(m,r)}\operatorname{exp}_2\left(-\epsilon\binom m{\le r}\right)\frac{\binom{(1-\mathrm{wt}(P))2^m}K}{\binom{2^m}K}\right].
\end{align*}
To complete the proof of Lemma 3.4 we therefore need the following claim.
\paragraph{Claim 3.5.}
\[\sum_{P\in\mathrm{RM}(m,r)}\operatorname{exp}_2\left(-\epsilon\binom m{\le r}\right)\frac{\binom{(1-\mathrm{wt}(P))2^m}K}{\binom{2^m}K}\le1+o(1).\]
Given Claim 3.5, we see that
\begin{align*}\mathbb E\left[\frac{|I_r(Z)|}{\operatorname{exp}_2\left(\binom m{\le r}\right)}\right]&=\operatorname{exp}_2(-K)\left[\sum_{P\in\mathrm{RM}(m,r)}\operatorname{exp}_2\left(-\epsilon\binom m{\le r}\right)\frac{\binom{(1-\mathrm{wt}(P))2^m}K}{\binom{2^m}K}\right]\\&\le\operatorname{exp}_2(-K)(1+o(1)),\end{align*}
which concludes the proof of Lemma 3.4. Next, we proceed to prove Claim 3.5.
\paragraph{Proof of Claim 3.5.} Let $u=\binom m{\le r}$ in the following. Recall that $K=(1-\epsilon)u$. Also,
\[\sum_{P\in\mathrm{RM}(m,r)}\operatorname{exp}_2(-\epsilon u)\frac{\binom{(1-\mathrm{wt}(P))2^m}K}{\binom{2^m}K}\le\sum_{P\in\mathrm{RM}(m,r)}\operatorname{exp}_2(-\epsilon u)(1-\mathrm{wt}(P))^{(1-\epsilon)u}.\]
Set $\delta=\left(1/\binom m{\le r}\right)^2$ and let $\mathrm{RM}_\delta(m,r)=\{P\in\mathrm{RM}(m,r):|\mathrm{wt}(P)-1/2|\ge\delta/2\}$, i.e., the set of polynomials with bias at least $\delta$. Now,
\begin{multline*}\sum_{P\in\mathrm{RM}_\delta(m,r)}\operatorname{exp}_2(-\epsilon u)(1-\mathrm{wt}(P))^{(1-\epsilon)u}\\\le2\sum_{P:\mathrm{wt}(P)\le1/2-\delta/2}\operatorname{exp}_2(-\epsilon u)(1-\mathrm{wt}(P))^{(1-\epsilon)u},\end{multline*}
since the term corresponding to a polynomial $P$ of weight greater than $1/2+\delta/2$ can be upper bounded by the term corresponding to the polynomial $1+P$ which has weight at most $1/2-\delta/2$.
Since the RHS of the above equation involves polynomials whose weights are in the interval $[1/2^r,1/2-\delta]$, we will split this interval into sub-intervals and analyse the contribution from each.
\begin{align*}\mathrm{Low}_i&:=[1/2^{i+1},1/2^i]&&\text{for }i=2,3,\ldots,r-1,\\\mathrm{Med}_i&:=[(1/2-1/2^i),(1/2-1/2^{i+1})]&&\text{for }i=2,3,\ldots,t=\log(1/\delta).\end{align*}
Let us use the following quantities to denote the number of polynomials with weights in the above intervals:
\[L_i:=|\{P\in\mathrm{RM}(m,r):\mathrm{wt}(P)\in\mathrm{Low}_i\}|,\qquad M_i:=|\{P\in\mathrm{RM}(m,r):\mathrm{wt}(P)\in\mathrm{Med}_i\}|.\]
Therefore,
\begin{align*}&\sum_{P\in\mathrm{RM}_\delta(m,r)}\operatorname{exp}_2(-\epsilon u)(1-\mathrm{wt}(P))^{(1-\epsilon)u}\\&\quad\le2\sum_{P:\mathrm{wt}(P)\le1/2-\delta/2}\operatorname{exp}_2(-\epsilon u)(1-\mathrm{wt}(P))^{(1-\epsilon)u}\\&\quad\le2\left(\sum_{i=2}^{r-1}L_i(1-1/2^{i+1})^k+\sum_{i=2}^tM_i(1/2+1/2^i)^k\right)\frac1{\operatorname{exp}_2\left(\binom m{\le r}-K\right)}.\end{align*}
By Claim 3.6 proved below using the weight distribution bounds of Sberlo and Shpilka (Theorem 2.1 and Theorem 2.2), we have that
\[\left(\sum_{i=2}^{r-1}L_i(1-1/2^{i+1})^k+\sum_{i=2}^tM_i(1/2+1/2^i)^k\right)\frac1{\operatorname{exp}_2\left(\binom m{\le r}-K\right)}\le O(\delta^2).\]
Hence,
\begin{align*}&\sum_{P\in\mathrm{RM}(m,r)}\operatorname{exp}_2(-\epsilon u)(1-\mathrm{wt}(P))^{(1-\epsilon)u}\\&=\sum_{P\in\mathrm{RM}_\delta(m,r)}\operatorname{exp}_2(-\epsilon u)(1-\mathrm{wt}(P))^{(1-\epsilon)u}\\&\qquad+\sum_{P\notin\mathrm{RM}_\delta(m,r)}\operatorname{exp}_2(-\epsilon u)(1-\mathrm{wt}(P))^{(1-\epsilon)u}\\&\le O(\delta^2)+\sum_{P\notin\mathrm{RM}_\delta(m,r)}\operatorname{exp}_2(-\epsilon u)(1/2+\delta/2)^{(1-\epsilon)u}\\&\le O(\delta^2)+(1+\delta)^{(1-\epsilon)u}\\&\le O(\delta^2)+\exp(\delta u)\\&\le1+O(1/u)=1+o(1).\end{align*}\qed
It remains to prove the following technical claim.
\paragraph{Claim 3.6.} 1. For all $i=2,\ldots,r-1$, we have
\[\frac{L_i(1-1/2^{i+1})^k}{\operatorname{exp}_2\left(\binom m{\le r}-k\right)}\le\delta^3.\]
2. For all $i=2,\ldots,t=\log(1/\delta)$, we have
\[\frac{M_i(1/2+1/2^i)^k}{\operatorname{exp}_2\left(\binom m{\le r}-k\right)}\le\delta^3.\]
\paragraph{Prooff of Claim 3.6(1).} Note that $(1-1/2^{i+1})^k\le\exp(-k/2^{i+1})\le\operatorname{exp}_2(-k/2^{i+1})$. Hence,
\[\frac{L_i(1-1/2^{i+1})^k}{\operatorname{exp}_2\left(\binom m{\le r}-k\right)}\le\frac{\mathrm{WtDist}_{m,r}(2^{-i})\operatorname{exp}_2(-k/2^{i+1})}{\operatorname{exp}_2\left(\binom m{\le r}-k\right)}.\]
Using Theorem 2.1 to bound $\mathrm{WtDist}_{m,r}(2^{-i})$, we get
\begin{align*}\frac{L_i(1-1/2^{i+1})^k}{\operatorname{exp}_2\left(\binom m{\le r}-k\right)}&\le\operatorname{exp}_2\left(O(m^4)+O(i\gamma^{i-1})\binom m{\le r}-k/2^{i+1}-\binom m{\le r}+k\right)\\&=\operatorname{exp}_2\left(O(m^4)-\binom m{\le r}(1-O(i\gamma^{i-1}))+k(1-2^{-i-1})\right),\end{align*}
where $\gamma=r/m$. If this $\gamma$ is a small enough absolute constant, we can see that the $O(i\gamma^{i-1})$ term is always at most $2^{-i-1}$. Using this and continuing the computation above, we get
\begin{align*}\frac{L_i(1-1/2^{i+1})^k}{\operatorname{exp}_2\left(\binom m{\le r}-k\right)}&\le\operatorname{exp}_2\left(O(m^4)-\left(\binom m{\le r}-k\right)(1-2^{-i-1})\right)\\&\le\operatorname{exp}_2\left(O(m^4)-\epsilon\binom m{\le r}(1-2^{-i-1})\right)\\&\le\operatorname{exp}_2\left(O(m^4)-\frac\epsilon2\binom m{\le r}\right)\\&\le\operatorname{exp}_2\left(-\Omega\left(\binom m{\le r}\right)\right)\le\delta^3\end{align*}
where for the second inequality above, we used the fact that $k\le K=(1-\epsilon)\binom m{\le r}$.\qed
\paragraph{Proof of Claim 3.6(2).} Note that
\begin{align*}(1/2+1/2^i)^k&=2^{-k}(1+1/2^{i-1})^k\\&\le2^{-k}\exp(k/2^{i-1})=2^{-k}\operatorname{exp}_2(k/2^{i-1}\cdot\log_2e)\\&=\operatorname{exp}_2(-k(1-2\log_2e\cdot2^{-i})).\end{align*}
Using the fact that $M_i\le\mathrm{WtDist}_{m,r}(1/2-1/2^{i+1})$, we get
\begin{align*}&\frac{\mathrm{WtDist}_{m,r}(1/2-1/2^{i+1})\operatorname{exp}_2(-k(1-2\log_2e\cdot2^{-i}))}{\operatorname{exp}_2\left(\binom m{\le r}-k\right)}\\&\quad\le\frac{\operatorname{exp}_2\left(O(m^4)+(1-2^{-O(\max\{\gamma^2i,\gamma\})})\binom m{\le r}-k(1-2\log_2e\cdot2^{-i}))\right)}{\operatorname{exp}_2\left(\binom m{\le r}-k\right)}\\&\quad=\operatorname{exp}_2\left(O(m^4)-\frac1{2^{O(\max\{\gamma^2i,\gamma\})}}\binom m{\le r}+k\cdot\frac1{2^{-(i-1-\log_2\log_2e)}}\right)\end{align*}
where $\gamma=r/m$, and we used Theorem 2.2 for the second inequality.
Note that as long as $\gamma$ is a small enough absolute constant, the $O(\max\{\gamma^2i,\gamma\})$ term above is at most $i-1-\log_2\log_2e$ for each $i\in\{2,\ldots,t\}$. Using this bound, we continue the above computation as follows.
\begin{align*}\frac{M_i(1/2+1/2^i)^k}{\operatorname{exp}_2\left(\binom m{\le r}-k\right)}&\le\operatorname{exp}_2\left(O(m^4)-2^{-O(\max\{\gamma^2i,\gamma\})}\left(\binom m{\le r}-k\right)\right)\\&\le\operatorname{exp}_2\left(O(m^4)-2^{-O(\max\{\gamma^2i,\gamma\})}\left(\binom m{\le r}-K\right)\right)\\&=\operatorname{exp}_2\left(O(m^4)-2^{-O(\max\{\gamma^2i,\gamma\})}\epsilon\binom m{\le r}\right)\\&\le\operatorname{exp}_2\left(O(m^4)-\Omega\left(\sqrt{\binom m{\le r}}\right)\right)\le\delta^3,\end{align*}
where for the last inequality we used the fact that $\binom m{\le r}\ge2^{\Omega(\gamma\log(1/\gamma)m)}$ and hence for small enough absolute constant $\gamma$, we have
\begin{align*}2^{-O(\max\{\gamma^2i,\gamma\})}\binom m{\le r}&\ge2^{-O(\gamma^2t)}\binom m{\le r}\\&\ge\operatorname{exp}_2(\Omega(\gamma\log(1/\gamma)m)-O(\gamma^2\log(1/\delta)))\\&=\operatorname{exp}_2(\Omega(\gamma\log(1/\gamma)m)-O(\gamma^3\log(1/\gamma)m))\\&=\operatorname{exp}_2(\Omega(m)).\end{align*}\qed
This completes the proof of Lemma 3.4.

\section*{Author Section 4: resilience under erasures and errors}
For parameters $m,r$, recall that $E(m,r)$ is the $\binom m{\le r}\times2^m$-matrix where the columns are indexed by elements in $\mathbb F_2^m$, and the $z$-th column is the column vector consisting of evaluations of all multilinear $m$-variate degree $\le r$ monomials on $z$. It is well-known that the matrix $E(m,r)$ is the generator matrix of the $\mathrm{RM}(m,r)$ code, and is also the parity-check matrix of the $\mathrm{RM}(m,m-r-1)$ code. Theorem 1.1 can be re-interpreted as a statement about random sub-matrices of $E(m,r)$ in the regime where $r\le\gamma m$ for a small absolute constant $\gamma>0$.
\paragraph{Corollary 4.1 (rank of random sub-matrices of $E(m,r)$).} Let $\gamma>0$ be an absolute constant and $m,r$ growing parameters with $r\le\gamma m$. Then, for any constant $\epsilon>0$, a random set of $K=(1-\epsilon)\binom m{\le r}$ columns of $E(m,r)$ are linearly independent with probability $1-o(1)$.
\begin{proof}A set of columns of $E(m,r)$ indexed by $Z=\{z_1,\ldots,z_K\}$ are linearly dependent if and only if $I_r(Z)$ has dimension strictly larger than $\binom m{\le r}-K$. Theorem 1.1 asserts that this happens with only $o(1)$ probability.\end{proof}
\subsection*{Author Section 4.1: channels under consideration}
The above corollary can be used to talk about the resilience of Reed-Muller codes under the erasure and error channels. We first define the precise definitions to be able to state the results accurately.
\paragraph{Definition 4.2 (binary erasure channel (BEC)).} The binary erasure channel with parameter $\alpha$, denoted by $\mathrm{BEC}_\alpha$, is the channel with input alphabet $\{0,1\}$ where the each binary symbol is ``erased'' (replaced by the `?' symbol) independently with probability $\alpha$.
A closely related model is where we have a fixed number of random erasures instead of each coordinate being erased with a fixed probability. We will refer to this as the $\mathrm{BEC}^*$ model although this isn't a channel in the traditional sense of altering each coordinate independently.
\paragraph{Definition 4.3 (capped binary erasure channel (BEC*)).} The capped binary erasure channel with parameter $K$, denoted by $\mathrm{BEC}_K^*$, refers to the channel with input alphabet $\{0,1\}^n$ that replaces a random subset of at most $K$ of the coordinates by the `?' symbol.
The Binary Symmetric Channel deals with errors or corruptions as opposed to erasures in the Binary Erasure Channel.
\paragraph{Definition 4.4 (binary symmetric channel (BSC)).} The binary symmetric channel with parameter $\alpha$, is the channel where the each binary symbol is ``flipped'' (that is, 0 changed to 1 and vice-versa) independently with probability $\alpha$.
In similar spirit to Definition 4.3, we define the capped binary symmetric channel.
\paragraph{Definition 4.5 (capped binary symmetric channel (BSC*)).} The capped binary symmetric channel with parameter $K$, denoted by $\mathrm{BSC}_K^*$, refers to the channel with input alphabet $\{0,1\}^n$ that ``flips'' a random subset of at most $K$ of the coordinates.
In most cases, resilience with respect to $\mathrm{BEC}_\alpha$ (or $\mathrm{BSC}_\alpha$) is the same as resilience with respect to $\mathrm{BEC}_K^*$ (or $\mathrm{BSC}_K^*$) for $K=\alpha n$, by standard concentration inequalities but this might require some subtlety when dealing with codes of rate very close to zero or close to one. For a concrete example, the code $\mathrm{RM}(m,m-1)$ has rate $R=1-1/2^m$ and is not resilient under $\mathrm{BEC}_\alpha$ for $\alpha=(1-R)/2=1/2^{m+1}$ (with constant probability we may have two coordinates erased, and this is not recoverable), but is vacuously resilient under $\mathrm{BEC}_K^*$ for $K=\alpha2^m$. To avoid these nuances, we will only be dealing with the setting of capped channels although this does not make much difference with most of our range of parameters.
\paragraph{Notion of ``capacity achieving'' codes with respect to $\mathrm{BEC}_\alpha$.} Shannon's seminar work [14] showed that, for any constant rate, the supremum of rates of random codes that is resilient to $\mathrm{BEC}_\alpha$ is exactly $R=1-\alpha$. That is, for any $\epsilon>0$, there are codes of rate $1-\alpha-\epsilon$ that can decode from $\mathrm{BEC}_\alpha$ with decoding error $1-o(1)$. However, when the rate is very close to zero or one, the situation becomes more nuanced due to the asymptotics involved.
Dealing specifically with Reed-Muller codes, let us consider the code $\mathrm{RM}(m,m-r-1)$ with $r\le m/2$. This has rate $R=1-\binom m{\le r}/2^m$ and can recover from a maximum of $\binom m{\le r}$ errors. Abbe, Shpilka and Wigderson [1] defined the notion of ``capacity achieving under BEC'' to mean that the code can decode (with $1-o(1)$ probability) from $(1-\epsilon)\binom m{\le r}$ erasures for any constant $\epsilon>0$. We refer the reader to [1, Section 1] and [13, Section 2.2] for a nuanced discussion on this.
\subsection*{Author Section 4.2: Reed-Muller codes under BEC*}
An immediate consequence of the above corollary is that $\mathrm{RM}(m,d)$ with $d\ge m(1-\gamma)$ for a small but absolute constant $\gamma>0$ achieves capacity (as defined above) under the Binary Erasure Channel.
\paragraph{Proof of Theorem 1.2.} The code $\mathrm{RM}(m,d)$ can recover from erasures on coordinates indexed by $Z\subseteq\mathbb F_2^m$ if and only if the corresponding columns of the parity check matrix are linearly independent. Since the parity-check matrix of $\mathrm{RM}(m,d)$ is $E(m,r)$ for $r=m-d-1$, we have from Corollary 4.1 that a random set of $K$ columns are linearly independent with probability $1-o(1)$.\qed
\subsection*{Author Section 4.3: Reed-Muller codes under BSC*}
The following theorem of Abbe, Shpilka and Wigderson [1] provides a method to derive the resilience of certain Reed-Muller codes under the BSC by using the resilience of appropriate Reed-Muller codes under the BEC. Subsequent work of Saptharishi, Shpilka and Volk also showed that these codes also have efficient decoding procedures to recover from random errors.
\paragraph{Theorem 4.6 ([1,12]).} For any growing parameters $m,r>0$, if the code $\mathrm{RM}(m,m-2r-2)$ can recover from a subset $Z\subseteq[2^m]$ of erasures, then the code $\mathrm{RM}(m,m-r)$ can efficiently recover from errors on the subset $Z$. In particular, if the code $\mathrm{RM}(m,m-2r-2)$ can recover $K$ random erasures with probability $1-o(1)$, then the code $\mathrm{RM}(m,m-r)$ can efficiently decode from $K$ random errors.
Applying the above theorem to Theorem 1.2 yields Corollary 1.3.
\paragraph{Corollary 4.7.} There exists a constant $\gamma_0>0$ such that for all $\gamma<\gamma_0$ and $\epsilon>0$ the following is true. Let $m,r$ be growing parameters with $r=\gamma m$. Then, the code $\mathrm{RM}(m,m-2r-2)$ can be efficiently decode from $K=(1-\epsilon)\binom m{\le r}$ random errors.
It is worth noting that the minimum distance of the code $\mathrm{RM}(m,m-2r-2)$ is merely $2^{2r+2}\ll\binom m{\le r}$ when $r=\gamma m$ for a small enough $\gamma>0$. Thus, the above corollary shows that high-degree (or high-rate) Reed-Muller codes are resilient to random errors well beyond their minimum distance, and efficiently so.

\begin{thebibliography}{15}
\bibitem{ref1} Emmanuel Abbe, Amir Shpilka, and Avi Wigderson. Reed-Muller codes for random erasures and errors. IEEE Trans. Inform. Theory, 61(10):5229–5252, 2015. (Preliminary version in 47th STOC, 2015). doi:10.1109/TIT.2015.2462817.
\bibitem{ref2} Ido Ben-Eliezer, Rani Hod, and Shachar Lovett. Random low-degree polynomials are hard to approximate. Comput. Complexity, 21(1):63–81, 2012. (Preliminary version in 13th RANDOM, 2009). eccc:2008/TR08-080. doi:10.1007/s00037-011-0020-6.
\bibitem{ref3} Aidao Chen, Anindya De, and Aravindan Vijayaraghavan. Learning a mixture of two subspaces over finite fields. In Vitaly Feldman, Katrina Ligett, and Sivan Sabato, editors, Algorithmic Learning Theory (ALT), volume 132 of Proceedings of Machine Learning Research, pages 481–504, 2021. arXiv:2010.02841.
\bibitem{ref4} Tali Kaufman, Shachar Lovett, and Ely Porat. Weight distribution and list-decoding size of Reed-Muller codes. IEEE Trans. Inform. Theory, 58(5):2689–2696, 2012. (Preliminary version in 1st ICS, 2010). doi:10.1109/TIT.2012.2184841.
\bibitem{ref5} Peter Keevash and Benny Sudakov. Set systems with restricted cross-intersections and the minimum rank of inclusion matrices. SIAM J. Discrete Math., 18(4):713–727, 2005. doi:10.1137/S0895480103434634.
\bibitem{ref6} Shrinivas Kudekar, Santhosh Kumar, Marco Mondelli, Henry D. Pfister, Eren Eren Şaşoğlu, and Rüdiger L. Urbanke. Reed-Muller codes achieve capacity on erasure channels. IEEE Trans. Inform. Theory, 63(7):4298–4316, 2017. (Preliminary version in 48th STOC, 2016). doi:10.1109/TIT.2017.2673829.
\bibitem{ref7} David E. Muller. Application of Boolean algebra to switching circuit design and to error detection. Trans. IRE Prof. Group Electron. Comput., 3(3):6–12, 1954. doi:10.1109/IREPGELC. 1954.6499441.
\bibitem{ref8} Zipei Nie and Anthony Y. Wang. Hilbert functions and the finite degree Zariski closure in finite field combinatorial geometry. J. Comb. Theory, Ser. A, 134:196–220, 2015. doi: 10.1016/j.jcta.2015.03.011.
\bibitem{ref9} Irving S. Reed. A class of multiple-error-correcting codes and the decoding scheme. Trans. IRE Prof. Group Inf. Theory, 4:38–49, 1954. doi:10.1109/TIT.1954.1057465.
\bibitem{ref10} Galen Reeves and Henry D. Pfister. Reed-Muller codes achieve capacity on BMS channels. (manuscript), 2021. arXiv:2110.14631.
\bibitem{ref11} Alex Samorodnitsky. An upper bound on $\ell$q norms of noisy functions. IEEE Trans. Inform. Theory, 66(2):742–748, 2020. doi:10.1109/TIT.2019.2944698.
\bibitem{ref12} Ramprasad Saptharishi, Amir Shpilka, and Ben Lee Volk. Efficiently decoding Reed-Muller codes from random errors. IEEE Trans. Inform. Theory, 63(4):1954–1960, 2017. (Preliminary version in 48th STOC, 2016). doi:10.1109/TIT.2017.2671410.
\bibitem{ref13} Ori Sberlo and Amir Shpilka. On the performance of Reed-Muller codes with respect to random errors and erasures. In Shuchi Chawla, editor, Proc. 31st Annual ACM-SIAM Symp. on Discrete Algorithms (SODA), pages 1357–1376, 2020. doi:10.1137/1.9781611975994.82.
\bibitem{ref14} Claude Elwood Shannon. A mathematical theory of communication. The Bell System Technical Journal, 27:379–423, 623–656, 1948. doi:10.1002/j.1538-7305.1948.tb01338.x.
\bibitem{ref15} Victor K.-W. Wei. Generalized Hamming weights for linear codes. IEEE Trans. Inform. Theory, 37(5):1412–1418, 1991. doi:10.1109/18.133259.
\end{thebibliography}
\end{document}
